\documentclass[UTF8]{ctexart}
\usepackage{geometry}
\usepackage[backend=biber, style=authoryear]{biblatex} % 示例配置
\addbibresource{references.bib} % 替换为你的参考文献文件
\geometry{margin=1.5cm, vmargin={0pt,1cm}}
\setlength{\topmargin}{-1cm}
\setlength{\paperheight}{29.7cm}
\setlength{\textheight}{25.3cm}

% useful packages.
\usepackage{amsfonts}
\usepackage{amsmath}
\usepackage{amssymb}
\usepackage{amsthm}
\usepackage{enumerate}
\usepackage{graphicx}
\usepackage{multicol}
\usepackage{fancyhdr}
\usepackage{layout}
\usepackage{listings}
\usepackage{xcolor}
\lstset{language=C++}
\setlength{\headheight}{13pt}  % 增加头部高度
\usepackage{float, caption}

% Listings setting for code formatting
\lstset{
    basicstyle=\ttfamily,
    keywordstyle=\color{blue},
    commentstyle=\color{gray},
    stringstyle=\color{red},
    breaklines=true,
    frame=single,
    numbers=left,
    numberstyle=\tiny,
    stepnumber=1,
    tabsize=4
}

% some common commands
\newcommand{\dif}{\mathrm{d}}
\newcommand{\avg}[1]{\left\langle #1 \right\rangle}
\newcommand{\difFrac}[2]{\frac{\dif #1}{\dif #2}}
\newcommand{\pdfFrac}[2]{\frac{\partial #1}{\partial #2}}
\newcommand{\OFL}{\mathrm{OFL}}
\newcommand{\UFL}{\mathrm{UFL}}
\newcommand{\fl}{\mathrm{fl}}
\newcommand{\op}{\odot}
\newcommand{\Eabs}{E_{\mathrm{abs}}}
\newcommand{\Erel}{E_{\mathrm{rel}}}

\begin{document}

\pagestyle{fancy}
\fancyhead{}
\lhead{陈梓锐}
\chead{22435036}
\rhead{\today}


%\begin{abstract}
 %   The abstract is not necessary for the theoretical homework, 
   % but for the programming project, 
    %you are encouraged to write one.      
%\end{abstract}

% ============================================
\subsection*{如何使用}
在code中build文件下使用make命令编译，编译成功后，在build文件夹下会生成可执行文件P,在test目录下的test.josn调整初始条件即可使用不同的测试用例。

在网格数为64时求解非常慢，，请耐心等待。

\subsection*{设计文档}
本次大作业要求做五点差分格式在规则或不规则区域，dirichlet或neumann边界条件下求解的偏微分方程的数值解。
所以只需要求解五点差分格式的矩阵方程组，将得到的解转化成图片即可。

首先，当输入的网格大小为m时，在所有求解点矩阵
$
\begin{pmatrix}
u_{11} & u_{12} & \cdots & u_{1m} \\
u_{21} & u_{22} & \cdots & u_{2m} \\
\vdots & \vdots & \ddots & \vdots \\
u_{m1} & u_{m2} & \cdots & u_{mm}
\end{pmatrix}
$之外，包围了一圈边界点，即$
\begin{pmatrix}
    u_{00} & u_{01} & \cdots & u_{0m} &u_{0,m+1}\\
    u_{10} & u_{11} & \cdots & u_{1m} &u_{1,m+1}\\
    \vdots & \vdots & \ddots & \vdots & \vdots \\
    u_{m0} & u_{m1} & \cdots & u_{mm} &u_{m,m+1}\\
    u_{m+1,0} & u_{m+1,1} & \cdots & u_{m+1,m} &u_{m+1,m+1}
\end{pmatrix}$
边界点外面再包一圈ghost point，但并不显式的存储在矩阵中，用于计算Neumann边界条件。

在构造矩阵方程组时，只需要考虑内部的点，所以系数矩阵A是一个m阶方阵。

\subsection*{程序缺陷}
算法在迭代五次后，残差和误差基本不变，多次检查后仍然无法确定问题所在。

\subsection*{Problem $exp(y+sin(x))$}
\subsubsection*{VCycle}
对于规则边界，Dirichlet边界条件，fullWeighting,linear。
\begin{enumerate}
\item 网格数为32。
    在第6次迭代时残差达到$10^{-6}$
    Residual norm at level 32: 0.000262642
    Residual norm at level 16: 4.34873e-05
    Residual norm at level 8: 8.95084e-06
    6 iteration, error = 1.01905


\item 网格数为64。
    在第6次迭代时残差达到$10^{-5}$
    Residual norm at level 64: 0.00178412
    Residual norm at level 32: 0.000330479
    Residual norm at level 16: 5.33846e-05
    Residual norm at level 8: 1.07992e-05
    6 iteration, error = 0.72165


\item 网格数为128。
    在第6次迭代时残差达到$10^{-5}$
    Residual norm at level 128: 0.011126
    Residual norm at level 64: 0.00208497
    Residual norm at level 32: 0.000374178
    Residual norm at level 16: 5.89672e-05
    Residual norm at level 8: 1.16673e-05
    6 iteration, error = 0.51057


\item 网格数为256。
    Residual norm at level 256: 0.0661172
    Residual norm at level 128: 0.0125497
    Residual norm at level 64: 0.00228404
    Residual norm at level 32: 0.000398829
    Residual norm at level 16: 6.16598e-05
    Residual norm at level 8: 1.20221e-05
    6 iteration, error = 0.361112

\end{enumerate}

规则边界，Dirichlet边界条件, fullWeighting, quadratic。
\item 网格数为32。
    在第6次迭代时残差达到$10^{-6}$
    Residual norm at level 32: 0.000243564
    Residual norm at level 16: 4.29247e-05
    Residual norm at level 8: 8.17586e-06
    6 iteration, error = 1.01905

\item 网格数为64。
    在第6次迭代时残差达到$10^{-6}$
    Residual norm at level 64: 0.00161502
    Residual norm at level 32: 0.000321223
    Residual norm at level 16: 5.24646e-05
    Residual norm at level 8: 9.46208e-06
    6 iteration, error = 0.721649

\item 网格数为128。
    在第6次迭代时残差达到$10^{-6}$
    Residual norm at level 128: 0.00931856
    Residual norm at level 64: 0.00193213
    Residual norm at level 32: 0.000358479
    Residual norm at level 16: 5.65874e-05
    Residual norm at level 8: 9.95138e-06
    6 iteration, error = 0.510569

\item 网格数为256。
    在第6次迭代时残差达到$10^{-5}$
    Residual norm at level 256: 0.0488913
    Residual norm at level 128: 0.0105387
    Residual norm at level 64: 0.00206747
    Residual norm at level 32: 0.000373161
    Residual norm at level 16: 5.80517e-05
    Residual norm at level 8: 1.01117e-05
    6 iteration, error = 0.361111

\end{enumerate}
说明二次插值的精度更高。

规则边界，Dirichlet边界条件, injection, linear。
\begin{enumerate}
\item 网格数为32。
    在第6次迭代时残差达到$10^{-5}$
    Residual norm at level 32: 0.00125484
    Residual norm at level 16: 0.000225143
    Residual norm at level 8: 4.71122e-05
    6 iteration, error = 1.01905
\item 网格数为64。
    在第6次迭代时残差达到$10^{-5}$
    Residual norm at level 64: 0.00827634
    Residual norm at level 32: 0.00167691
    Residual norm at level 16: 0.000276912
    Residual norm at level 8: 5.52438e-05
    6 iteration, error = 0.721652
\item 网格数为128。
    在第6次迭代时残差达到$10^{-5}$
    Residual norm at level 128: 0.0500838
    Residual norm at level 64: 0.01032
    Residual norm at level 32: 0.00191325
    Residual norm at level 16: 0.000303276
    Residual norm at level 8: 5.87188e-05
    6 iteration, error = 0.510574
\item 网格数为256。
    在第6次迭代时残差达到$10^{-5}$
    Residual norm at level 256: 0.285242
    Residual norm at level 128: 0.0599207
    Residual norm at level 64: 0.0113986
    Residual norm at level 32: 0.00202851
    Residual norm at level 16: 0.000314789
    Residual norm at level 8: 6.00298e-05
    6 iteration, error = 0.361119
\end{enumerate}
说明fullweighting松弛的精度更高。

规则边界，Dirichlet边界条件, injection, quadratic。
\begin{enumerate}
\item 网格数为32。
    在第6次迭代时残差达到$10^{-5}$
    Residual norm at level 32: 0.00109206
    Residual norm at level 16: 0.000211448
    Residual norm at level 8: 4.30871e-05
    6 iteration, error = 1.01905
\item 网格数为64。
    在第6次迭代时残差达到$10^{-5}$
    Residual norm at level 64: 0.00664948
    Residual norm at level 32: 0.0014904
    Residual norm at level 16: 0.000257301
    Residual norm at level 8: 4.89599e-05
    6 iteration, error = 0.721652
\item 网格数为128。
    在第6次迭代时残差达到$10^{-5}$
    Residual norm at level 128: 0.0355607
    Residual norm at level 64: 0.00829714
    Residual norm at level 32: 0.0016712
    Residual norm at level 16: 0.000275795
    Residual norm at level 8: 5.10984e-05
    6 iteration, error = 0.510573
\item 网格数为256。
    在第6次迭代时残差达到$10^{-5}$
    Residual norm at level 256: 0.174971
    Residual norm at level 128: 0.0422026
    Residual norm at level 64: 0.00898611
    Residual norm at level 32: 0.00173936
    Residual norm at level 16: 0.000282247
    Residual norm at level 8: 5.17974e-05
    6 iteration, error = 0.361116
\end{enumerate}


\subsection*{误差分析}
观察误差发现每次h增加一倍，误差减小一倍，说明格式是$O(h^2)$的。









% ===============================================

\end{document}